
\SGARefTarget{sga3:i:subsection:1:8}
\subsection*{1.8. Constant objects}

Let \(\mathscr C\) be a category in which direct sums and fiber products
exist and in which direct sums commute with base change (for example, the
category of schemes). For every set \(E\) and every object \(S\) of
\(\mathscr C\), let \(E_S\) be the direct sum of a family
\((S_i)_{i\in E}\) of objects of \(\mathscr C\) all isomorphic to \(S\).
This object is characterized by the formula
\[
  \operatorname{Hom}(E_S,T)
  =\operatorname{Hom}\bigl(E,\operatorname{Hom}(S,T)\bigr),
  \qquad T\in\operatorname{Ob}(\mathscr C),
\]
where the second \(\operatorname{Hom}\) is taken in the category of sets.

The object \(E_S\) has a canonical projection to \(S\), so that
\(E\mapsto E_S\) is in fact a functor from \(\mathbf{Set}\) to
\(\mathscr C_{/S}\). Suppose that, with \(\varnothing\) denoting an initial
object of \(\mathscr C\), the diagram

\begin{center}
  \begin{tikzcd}[column sep=4.2em,row sep=3.0em]
  \varnothing
    \arrow[r]
    \arrow[d]
  &
  S
    \arrow[d]
  \\
  S
    \arrow[r]
  &
  S\amalg S
\end{tikzcd}

\end{center}
is cartesian (this is the case for the category of schemes). Then the
functor \(E\mapsto E_S\) commutes with finite inverse limits. In
particular,
\[
  E_S\mathbin{\times_S}F_S=(E\times F)_S.
\]
If \(S'\to S\) is an arrow of \(\mathscr C\), then
\[
  E_{S'}=(E_S)_{S'}.
\]
In particular, if \(\mathscr C\) has a final object \(e\), then
\[
  E_S=(E_e)_S.
\]

\par\medskip\noindent
\emph{Addition} (after the Gille--Polo re-edition (2011), Exposé~I, 1.8;
our formulation). The formula \(E_{S'}=(E_S)_{S'}\) is the commutation of
direct sums with base change. The functor \(E\mapsto E_S\) commutes with
finite products even when the square above is not cartesian: the one-point
set goes to \(S\), the final object of \(\mathscr C_{/S}\), and
\(E_S\times_SF_S\), the base change of \(E_S=\coprod_{i\in E}S_i\) by
\(F_S\to S\), is
\(\coprod_{i\in E}F_S=\coprod_{(i,j)\in E\times F}S=(E\times F)_S\). This
suffices to carry structures defined by finite products (2.2). When the
square is cartesian, \(E\mapsto E_S\) also commutes with fibre products,
hence with finite inverse limits. First, an initial object \(\varnothing\)
is strict: if \(X\to\varnothing\) is a morphism, the base change by it of
the empty sum \(\varnothing\) of \(\mathscr C_{/\varnothing}\) is \(X\), and
must be the empty sum of \(\mathscr C_{/X}\), so \(X\simeq\varnothing\).
Now let \(u:E\to G\), \(v:F\to G\) be maps; base change gives
\(E_S\times_{G_S}F_S=\coprod_{(i,j)\in E\times F}S_i\times_{G_S}S_j\). If
\(u(i)\ne v(j)\), a morphism \(G_S\to S\amalg S\) sending \(S_{u(i)}\) and
\(S_{v(j)}\) to different summands maps \(S_i\times_{G_S}S_j\) to
\(S\times_{S\amalg S}S\simeq\varnothing\), so
\(S_i\times_{G_S}S_j\simeq\varnothing\). If \(u(i)=v(j)=k\), base change of
\(G_S=S_k\amalg R\) by \(S_k\to G_S\) gives
\(S_k=(S_k\times_{G_S}S_k)\amalg(R\times_{G_S}S_k)\), the second term being
initial by the same argument; so \(S_k\to G_S\) is a monomorphism and
\(S_i\times_{G_S}S_j\simeq S\). Hence
\(E_S\times_{G_S}F_S\simeq\coprod_{E\times_GF}S=(E\times_GF)_S\).
\par\medskip

An object of the form \(E_S\) is called a \emph{constant object}. Note that
there is a morphism, functorial in \(E\),
\[
  E\longrightarrow\Gamma(E_S/S),
\]
which associates with each \(i\in E\) the section of \(E_S\) over \(S\)
defined by the isomorphism from \(S\) onto \(S_i\); it is a monomorphism if
the square above is cartesian and \(S\) is not an initial
object.\footnote{Translator's note: the 1970 text calls this morphism a
monomorphism without condition; the Gille--Polo re-edition (2011),
Exposé~I, adds conditions, and some are needed. If \(S\) is an initial
object, so is \(E_S\), and \(\Gamma(E_S/S)\) has one element. In the
category with two objects \(A\), \(B\) and one arrow \(A\to B\) besides the
identities, which satisfies the hypotheses of 1.8 but in which the square is
not cartesian, \(B\amalg B=B\). Under the two conditions, if \(i\ne j\) gave
the same section, \(S\) would map to \(S_i\times_{E_S}S_j\), which maps to
the initial object \(S\times_{S\amalg S}S\) through a morphism
\(E_S\to S\amalg S\) sending \(S_i\) and \(S_j\) to different summands; so
\(S\) would be initial, the initial object being strict (Addition above).}

If \(\mathscr C\) is the category of schemes, then
\(\Gamma(E_S/S)\) is identified with the locally constant maps from the
topological space \(S\) to the set \(E\), the preceding map associating with
each element of \(E\) the corresponding constant map.\footnote{Translator's note: the Gille--Polo re-edition (2011), Exposé~I, adds that
\(E_S\) is separated over \(S\): the diagonal
\(E_S\to E_S\times_SE_S=(E\times E)_S\) is the inclusion of the summands
indexed by the diagonal of \(E\times E\), an open and closed immersion.} It follows that
\(E_S\) may equivalently be defined as representing the functor that assigns
to every \(S'\) over \(S\) the set of locally constant functions from the
topological space \(S'\) to \(E\).

\section*{\SGAInlineRefTarget{sga3:i:section:2}2. Algebraic structures}
\addcontentsline{toc}{section}{2. Algebraic structures}

Given a species of algebraic structure in the category of sets, we shall
extend it to the category \(\mathscr C\). We first treat the example of
groups.

\subsection*{\SGAInlineRefTarget{sga3:i:subsection:2-1}2.1. Group structures}

We retain the notation of the preceding paragraph.

\medskip
\noindent\SGAInlineRefTarget{sga3:i:definition:2-1-1}\textbf{Definition 2.1.1.}
Let \(\mathbf G\in\operatorname{Ob}(\widehat{\mathscr C})\). A
\emph{\(\widehat{\mathscr C}\)-group structure} on \(\mathbf G\) is the
datum, for every \(S\in\operatorname{Ob}(\mathscr C)\), of a group structure
on the set \(\mathbf G(S)\), such that for every arrow
\(f:S'\to S''\) of \(\mathscr C\), the map
\[
  \mathbf G(f):\mathbf G(S'')\longrightarrow\mathbf G(S')
\]
is a group homomorphism. If \(\mathbf G\) and \(\mathbf H\) are two
\(\widehat{\mathscr C}\)-groups, a \emph{morphism of
\(\widehat{\mathscr C}\)-groups} from \(\mathbf G\) to \(\mathbf H\) is a
morphism \(u\in\operatorname{Hom}(\mathbf G,\mathbf H)\) such that
\(u(S):\mathbf G(S)\to\mathbf H(S)\) is a group homomorphism for every
\(S\in\operatorname{Ob}(\mathscr C)\).

We denote by
\(\operatorname{Hom}_{\widehat{\mathscr C}\text{-}\mathrm{Gr}}
  (\mathbf G,\mathbf H)\)
the set of these morphisms and by
\((\widehat{\mathscr C}\text{-}\mathrm{Gr})\) their category.

\medskip
\noindent\textbf{Examples.}
For \(\mathbf E\in\operatorname{Ob}(\widehat{\mathscr C})\), the object
\(\underline{\operatorname{Aut}}(\mathbf E)\) has an evident
\(\widehat{\mathscr C}\)-group structure. The final object
\(\underline{\mathbf e}\) has a unique
\(\widehat{\mathscr C}\)-group structure, making it a final object of
\((\widehat{\mathscr C}\text{-}\mathrm{Gr})\).

For every \(S\in\operatorname{Ob}(\mathscr C)\), let
\(e_{\mathbf G}(S)\) be the identity element of \(\mathbf G(S)\). The family
of elements \(e_{\mathbf G}(S)\) defines
\[
  e_{\mathbf G}\in\Gamma(\mathbf G)
  =\operatorname{Hom}(\underline{\mathbf e},\mathbf G).
\]
It is a morphism of \(\widehat{\mathscr C}\)-groups
\(\underline{\mathbf e}\to\mathbf G\), called the \emph{unit section} of
\(\mathbf G\).

To give a \(\widehat{\mathscr C}\)-group structure on \(\mathbf G\) is
equivalently to give a composition law on \(\mathbf G\), that is, a
\(\widehat{\mathscr C}\)-morphism
\[
  \pi_{\mathbf G}:\mathbf G\times\mathbf G\longrightarrow\mathbf G
\]
such that \(\pi_{\mathbf G}(S)\) equips \(\mathbf G(S)\) with a group
structure for every \(S\in\operatorname{Ob}(\mathscr C)\).

Similarly, \(f:\mathbf G\to\mathbf H\) is a morphism of
\(\widehat{\mathscr C}\)-groups if and only if the following diagram is
commutative:

\begin{center}
  \begin{tikzcd}[column sep=4.5em,row sep=3.0em]
  \mathbf G\times\mathbf G
    \arrow[r,"\pi_{\mathbf G}"]
    \arrow[d,"{(f,f)}"']
  &
  \mathbf G
    \arrow[d,"f"]
  \\
  \mathbf H\times\mathbf H
    \arrow[r,"\pi_{\mathbf H}"]
  &
  \mathbf H
  \qquad .
\end{tikzcd}

\end{center}

A subobject \(\mathbf H\) of \(\mathbf G\) such that \(\mathbf H(S)\) is a
subgroup of \(\mathbf G(S)\) for every \(S\) has the evident
\(\widehat{\mathscr C}\)-group structure induced by that of \(\mathbf G\).
It is the unique structure for which the monomorphism
\(\mathbf H\to\mathbf G\) is a morphism of
\(\widehat{\mathscr C}\)-groups. With this structure, \(\mathbf H\) is
called a \emph{sub-\(\widehat{\mathscr C}\)-group} of \(\mathbf G\).

If \(\mathbf G\) and \(\mathbf H\) are two
\(\widehat{\mathscr C}\)-groups, then \(\mathbf G\times\mathbf H\) has the
evident product structure: for every \(S\), the set
\(\mathbf G(S)\times\mathbf H(S)\) is given the product of the prescribed
group structures. The resulting object is called the
\emph{product \(\widehat{\mathscr C}\)-group} of \(\mathbf G\) and
\(\mathbf H\); it is their product in the category of
\(\widehat{\mathscr C}\)-groups.

If \(\mathbf G\) is a \(\widehat{\mathscr C}\)-group, then
\(\mathbf G_S\) is a \(\widehat{\mathscr C_{/S}}\)-group for every
\(S\in\operatorname{Ob}(\mathscr C)\). For two
\(\widehat{\mathscr C}\)-groups \(\mathbf G\) and \(\mathbf H\), define
the object
\(\underline{\operatorname{Hom}}_{
  \widehat{\mathscr C}\text{-}\mathrm{Gr}}(\mathbf G,\mathbf H)\)
of \(\widehat{\mathscr C}\) by
\[
  \underline{\operatorname{Hom}}_{
    \widehat{\mathscr C}\text{-}\mathrm{Gr}}(\mathbf G,\mathbf H)(S)
  =
  \operatorname{Hom}_{
    \widehat{\mathscr C_{/S}}\text{-}\mathrm{Gr}}
    (\mathbf G_S,\mathbf H_S).
\]
Notice that
\(\underline{\operatorname{Hom}}_{
  \widehat{\mathscr C}\text{-}\mathrm{Gr}}\)
is not in general a \(\widehat{\mathscr C}\)-group, still less the internal
\(\underline{\operatorname{Hom}}\) object in the category
\((\widehat{\mathscr C}\text{-}\mathrm{Gr})\).

One similarly defines
\[
  \underline{\operatorname{Isom}}_{
    \widehat{\mathscr C}\text{-}\mathrm{Gr}}(\mathbf G,\mathbf H),
  \qquad
  \underline{\operatorname{End}}_{
    \widehat{\mathscr C}\text{-}\mathrm{Gr}}(\mathbf G),
  \qquad
  \underline{\operatorname{Aut}}_{
    \widehat{\mathscr C}\text{-}\mathrm{Gr}}(\mathbf G).
\]

\medskip
\noindent\SGAInlineRefTarget{sga3:i:definition:2-1-2}\textbf{Definition 2.1.2.}
Let \(G\in\operatorname{Ob}(\mathscr C)\). A
\emph{\(\mathscr C\)-group structure} on \(G\) is a
\(\widehat{\mathscr C}\)-group structure on
\(\mathbf h_G\in\operatorname{Ob}(\widehat{\mathscr C})\). A morphism from
the \(\mathscr C\)-group \(G\) to the \(\mathscr C\)-group \(H\) is an
element
\[
  u\in\operatorname{Hom}(G,H)
  \simeq\operatorname{Hom}(\mathbf h_G,\mathbf h_H)
\]
that defines a morphism of \(\widehat{\mathscr C}\)-groups from
\(\mathbf h_G\) to \(\mathbf h_H\).

We denote the category of \(\mathscr C\)-groups by
\((\mathscr C\text{-}\mathrm{Gr})\). There is a cartesian square in
\(\mathbf{Cat}\):

\begin{center}
  \begin{tikzcd}[column sep=5.0em,row sep=3.2em]
  (\mathscr C\text{-}\mathrm{Gr.})
    \arrow[r]
    \arrow[d]
  &
  (\widehat{\mathscr C}\text{-}\mathrm{Gr.})
    \arrow[d]
  \\
  \mathscr C
    \arrow[r,"\mathbf h"]
  &
  \widehat{\mathscr C}
  \qquad .
\end{tikzcd}

\end{center}

All the preceding definitions and constructions therefore carry over at
once to \((\mathscr C\text{-}\mathrm{Gr})\) whenever the functors they use
(products, internal \(\underline{\operatorname{Hom}}\) objects, and so on)
are representable. They also apply to the categories \(\mathscr C_{/S}\).
In that case we write
\(\underline{\operatorname{Hom}}_{S\text{-}\mathrm{Gr}}\) for
\(\underline{\operatorname{Hom}}_{\mathscr C_{/S}\text{-}\mathrm{Gr}}\),
and similarly for the other objects.

\subsection*{\SGAInlineRefTarget{sga3:i:subsection:2-2}2.2.}

More generally, suppose \((T)\) is a species of structure on \(n\) base sets
defined by finite inverse limits---for example, by commutative diagrams
constructed from cartesian products, as in the structures of a monoid, a
group, a set with operators, a module over a ring, or a Lie algebra over a
ring. The preceding construction defines the notion of a ``structure of
species \((T)\)'' on \(n\) objects
\(\mathbf F_1,\ldots,\mathbf F_n\) of \(\widehat{\mathscr C}\). Such a
structure consists, for every \(S\in\operatorname{Ob}(\mathscr C)\), of a
structure of species \((T)\) on the sets
\(\mathbf F_1(S),\ldots,\mathbf F_n(S)\), in such a way that for every arrow
\(S'\to S''\) of \(\mathscr C\), the family of maps
\[
  \bigl(\mathbf F_i(S'')\bigr)_i
  \longrightarrow
  \bigl(\mathbf F_i(S')\bigr)_i
\]
is a polyhomomorphism for the species \((T)\).

Morphisms of structures of species \((T)\) are defined similarly, giving a
category
\[
  (\widehat{\mathscr C}\times\cdots\times\widehat{\mathscr C})^{(T)}.
\]
The fully faithful functor \(\mathbf h\times\cdots\times\mathbf h\) then
defines, by inverse image, the category
\[
  (\mathscr C\times\cdots\times\mathscr C)^{(T)}.
\]
Since it commutes with inverse limits, all the functorial properties,
notions, and notation introduced in \(\widehat{\mathscr C}\) carry over to
this category.

Now suppose fiber products exist in \(\mathscr C\), and let \((T)\) be a
species of algebraic structure specified by certain morphisms between
cartesian products, subject to axioms consisting of commutative diagrams
formed with those morphisms. A structure of species \((T)\) on a family of
objects of \(\mathscr C\) is then defined by corresponding morphisms between
cartesian products satisfying the corresponding commutativity conditions.
Consequently, if \(\mathscr C\) and \(\mathscr C'\) are categories with
products and \(f:\mathscr C\to\mathscr C'\) is a functor commuting with
products, then a structure of species \((T)\) on a family
\((F_i)\) of objects of \(\mathscr C\) induces one on the family
\((f(F_i))\) in \(\mathscr C'\). Thus every \(\mathscr C\)-group becomes a
\(\mathscr C'\)-group, every pair consisting of a \(\mathscr C\)-ring and a
module over it becomes the analogous pair in \(\mathscr C'\), and so on.

In particular, let \(\mathscr C\) satisfy the conditions of \SGARefLink{sga3:i:subsection:1:8}{1.8}. The functor
\(E\mapsto E_S\) defined there commutes with finite inverse
limits.\footnote{Translator's note: the 1970 text has ``commute aux limites
projectives'', without ``finite''; infinite inverse limits are not preserved
in general (over a field, an infinite product of two-point sets gives a
constant scheme, with a discrete underlying space, while the product of the
corresponding constant schemes has a Cantor set as underlying space). The
Gille--Polo re-edition (2011), Exposé~I, adds ``finite'', and notes that
the cartesian condition of 1.8 is included among the conditions. For
carrying groups, rings and other structures defined by finite products,
that condition is not needed (Addition in 1.8).} It
therefore carries groups to \(S\)-groups (that is,
\(\mathscr C_{/S}\)-groups), rings to \(S\)-rings, and so forth.

\medskip
\noindent\textbf{Remark.}
It is worth observing that the preceding construction, when applied to
\(\widehat{\mathscr C}\), gives back precisely the notions already defined
there. In other words, to give an object of \(\widehat{\mathscr C}\) a
structure of species \((T)\) when it is viewed as a functor on
\(\mathscr C\) is the same as to give the representable functor on
\(\widehat{\mathscr C}\) defined by that object a structure of species
\((T)\).\footnote{Translator's note: the Gille--Polo re-edition (2011), Exposé~I, spells this out for
group structures. In our formulation: an object \(\mathbf G\) of
\(\widehat{\mathscr C}\) defines the functor
\(\mathbf F\mapsto\operatorname{Hom}(\mathbf F,\mathbf G)\) on
\(\widehat{\mathscr C}\), whose restriction to \(\mathscr C\) is
\(\mathbf G\) itself (1.1); a group structure on the former restricts to a
\(\widehat{\mathscr C}\)-group structure on \(\mathbf G\). Conversely, a
\(\widehat{\mathscr C}\)-group structure on \(\mathbf G\) is a
multiplication \(\mathbf G\times\mathbf G\to\mathbf G\) satisfying the group
axioms, and composition with it makes each
\(\operatorname{Hom}(\mathbf F,\mathbf G)\) a group, functorially in
\(\mathbf F\); the two constructions are inverse to each other.}

We shall treat two further special cases of this construction: structures
with operator groups and modules.

